\section{Introduction}\label{sec:introduction}

\subsection{Context: urban traffic and the external-influence problem}
\label{sec:context}

The widespread adoption of satellite-navigation applications such as Google
Maps, Waze, and TomTom has structurally altered how vehicles use urban road
networks. These systems no longer merely report a single shortest path:
they reroute users in real time around observed congestion, redistributing
traffic across alternative arteries on time-scales of minutes to hours.
The consequence is that any traffic observation aggregating across modern
drivers reflects a \emph{mixture} of two distinct routing rules: the
intrinsic preferences a driver would express in the absence of guidance
(habit, scenic preference, fuel-cost optimisation, familiarity), and the
external recommendations of a sat-navigation system reacting to live
network state.

This mixture creates a problem for downstream traffic models. Demand
estimation, route-choice inference, transport-policy evaluation, and
infrastructure-investment decisions all rely on revealed-preference data
that is assumed to encode \emph{driver} preferences. When a sizeable
fraction of observed trips are guided by an external recommender that
itself responds to congestion (and thereby to other guided trips), the
revealed data encodes a partially synthetic, congestion-coupled signal
that the underlying driver population would not have produced unaided.
Identifying and filtering this externally-influenced traffic, or at minimum
quantifying its presence, has therefore become a methodological
prerequisite for any inference about driver behaviour in modern urban
environments.

\subsection{Inverse Markov-chain methods and the Morimura framework}
\label{sec:intro-morimura}

A long line of work models route choice as a Markov chain over road
network edges, in which each driver, conditional on their current edge,
selects a next edge with probabilities determined by local and global
features of the network. \citet{morimura2013} introduced an
\emph{inverse} formulation that fits the parameters of such a chain
from only partial observations of the network --- specifically, from
vehicle counts at a small subset of observation stations $\Xo$ and
optionally from pairwise hitting rates between those stations. Their
softmax-parametric chain admits an analytic gradient of the regularised
log-ratio loss, enabling reliable fitting on networks of several
thousand edges with $20\text{--}90\%$ observation coverage.

The Morimura framework is well validated on synthetic data and was
applied by its authors to one real-world traffic case study in Nairobi.
It has not, however, been adapted to the specific problem of detecting
externally-influenced routing as distinct from intrinsic preferences:
the original work treats the inverse problem as one of \emph{recovering}
a single chain, not of \emph{distinguishing} two coexisting routing
regimes within the same network.

\subsection{Contributions}
\label{sec:contributions}

This dissertation extends the Morimura inverse-Markov-chain method into a
three-tier detection pipeline for externally-influenced routing, validates
each tier honestly, and quantifies where the framework's limits actually
lie. Specifically:

\begin{enumerate}
\item \textbf{Faithful synthetic reproduction of the Morimura $\S 6.1$
benchmark.} We re-implement the inverter, derive and
finite-difference-verify the analytic gradients (including the
omega-global gradient extensions for shared destination and edge features
that the paper states but does not derive in full), and reproduce the
proposed-method RMAE curve of Fig.~2A in \cite{morimura2013} closely for
observation fractions $\Xo/n \ge 0.10$. The very-low-observation point
($|\Xo|=5$) and the NWKR comparator are documented deviations
(\cref{sec:phase1-synthetic}). This locks in the upstream implementation
as a faithful operational reference.

\item \textbf{A SUMO regime detector under partial observation with real
edge features.} We construct a bigger-bbox SUMO scenario in which an
identical population of vehicles is simulated twice --- once with no
sat-navigation rerouting (static unguided routing) and once with full
sat-navigation rerouting --- fit a Morimura inverter to each regime's
partial observations, and feed the fitted chains to a per-trajectory
log-likelihood-ratio detector. On the representative seed the detector
reaches $\AUC = 0.644$ at $T=20$ against a full-observation empirical
$1$-step reference of $0.685$, rising to $\AUC = 0.708$ at $T=40$ against
a reference of $0.801$ as the scoring window widens.

\item \textbf{Per-row contrast diagnostics separating chain recovery from
aggregate AUC.} Following the observation that detection AUC can improve
through aggregate chain-drift rather than correct per-row recovery, we
report a contrast-correlation diagnostic at branching states. Real edge
features robustly improve \emph{per-row chain recovery} across a five-seed
sweep (stacked Pearson $+0.050 \pm 0.023$, $p \approx 0.008$, positive on
all five seeds), whereas the aggregate detection-AUC lift is
seed-dependent rather than robust ($+0.024 \pm 0.026$, reversing on one
of five seeds). Features are best described as a chain-recovery lever, not
an aggregate-detection lever.

\item \textbf{An empirical higher-order memory diagnostic.} A natural
hypothesis is that the residual gap is missing $2$-step memory. We test
this with a fitting-free empirical $2$-step detector --- counting
$(x_{t-1}, x_t)\to x_{t+1}$ triples on the same trajectories --- and find
it does \emph{not} improve over the $1$-step empirical reference
($\AUC_{2\text{-step}} = 0.666 < \AUC_{1\text{-step}} = 0.685$). This
weakens the case for implementing a parametric $2$-step extension as the
immediate next step, but does not rule out regularised, feature-based,
destination-aware, or hierarchical higher-order models.

\item \textbf{A real-world stress test with OD-rebalancing on the
Xuancheng AVI dataset.} We port the pipeline to $\sim 10.5$ million
reconstructed vehicle trips over April~2023 \cite{ma2026xuancheng}. Because
no real dataset comes with sat-navigation adoption labels, we use a
time-of-day behavioural proxy (rush-hour vs.\ off-peak) and introduce an
OD-rebalancing step that subsamples each regime to share the same
$(\mathrm{origin\ zone},\mathrm{destination\ zone})$ distribution. This
controls the population-mix confound and shows that most of the apparent
proxy contrast is OD-mix, not routing choice (\cref{sec:phase5-xuancheng}).
A single CLI runs all three tiers via a \code{--dataset \{sumo,xuancheng\}}
flag; the repository structure is documented in \cref{app:repo}.
\end{enumerate}

\subsection{Roadmap}
\label{sec:roadmap}

\cref{sec:background} situates the work in relevant prior literature.
\cref{sec:methodology} presents the mathematical framework and the three
extensions (real features, $2$-step diagnostic, OD-rebalancing) in detail.
\cref{sec:phase1-synthetic,sec:phase4-sumo,sec:phase5-xuancheng}
present the three validation tiers in turn, each one with its own
honest discussion of what does and does not transfer. \cref{sec:discussion}
draws cross-tier conclusions and limitations. \cref{sec:reflection}
addresses project management and personal-development reflection per the
module's \emph{Reflection} learning outcome \cite{markscheme}.
